\documentclass[../../main.tex]{subfiles}
\graphicspath{{\subfix{../../image/}}}

\begin{document}

\section{附加课题}\label{section:ch05-补充主题}

\subsection{Poincaré 不等式}\label{subsection:ch05-边界情形}

\begin{theorem}[Poincaré 不等式]\label{theorem:ch05-Poincare}
设 $U$ 是$\mathbb{R}^{n}$的具有$C^1$边界$\partial U$的有界连通开子集,$1\leqslant p\leqslant\infty$。则存在只依赖于 $n$、$p$ 与 $U$的常数 $C$，使对每个函数 $u\in W^{1,p}(U)$ 有
\begin{align}
\|u-(u)_{U}\|_{L^{p}(U)}\leqslant C\|Du\|_{L^{p}(U)}.
\end{align}
其中$(u)_{U}=\fint_{U}u\,\mathrm{d}y$，即 $u$ 在 $U$ 上的平均。
\end{theorem}
\begin{remark}
此式的意义在于右端只出现 $u$ 的梯度。
\end{remark}
\begin{proof}

反证。若所述估计不成立，则对每个整数 $k=1,2,\cdots$ 存在函数 $u_{k}\in W^{1,p}(U)$ 满足
\begin{align}\label{eq:iojiowjwoioioii}
\|u_{k}-(u_{k})_{U}\|_{L^{p}(U)}>k\|Du_{k}\|_{L^{p}(U)}.
\end{align}
我们归一化，定义
\begin{align*}
v_{k}\triangleq\frac{u_{k}-(u_{k})_{U}}{\|u_{k}-(u_{k})_{U}\|_{L^{p}(U)}}\qquad (k=1,2,\cdots).
\end{align*}
则
\begin{align*}
(v_{k})_{U}=0,\qquad \|v_{k}\|_{L^{p}(U)}=1;
\end{align*}
且上述不等式\eqref{eq:iojiowjwoioioii}蕴含
\begin{align}\label{eq::jiojwjjwiwIJIWD}
\|Dv_{k}\|_{L^{p}(U)}<\frac{1}{k}\qquad (k=1,2,\cdots).
\end{align}
特别地，函数 $\{v_{k}\}_{k=1}^{\infty}$ 在 $W^{1,p}(U)$ 中有界。

由\hyperref[theorem:Rellich--Kondrachov紧性定理]{Rellich--Kondrachov紧性定理}，存在子序列 $\{v_{k_{j}}\}_{j=1}^{\infty}\subset\{v_{k}\}_{k=1}^{\infty}$ 及函数 $v\in L^{p}(U)$ 使
\begin{align}
v_{k_{j}}\to v\quad\text{于}\ L^{p}(U).
\end{align}
由归一化式得
\begin{align*}
(v)_{U}=0,\qquad \|v\|_{L^{p}(U)}=1.
\end{align*}
另一方面，\eqref{eq::jiojwjjwiwIJIWD}式蕴含对每个 $i=1,\cdots,n$ 及 $\phi\in C^{\infty}_{c}(U)$，
\begin{align*}
\int_{U}v\phi_{x_{i}}\,\mathrm{d}x
=\lim_{k_{j}\to\infty}\int_{U}v_{k_{j}}\phi_{x_{i}}\,\mathrm{d}x
=-\lim_{k_{j}\to\infty}\int_{U}v_{k_{j},x_{i}}\phi\,\mathrm{d}x=0.
\end{align*}
因而 $v\in W^{1,p}(U)$，且 $Dv(x)=0$ a.e.$\,x\in U$。于是由于 $U$ 连通和\cref{proposition:ch05--examoi-11}知，$v$ 为常数。然而此结论与 $\|v\|_{L^{p}(U)}=1$、$(v)_{U}=0$ 矛盾!
\end{proof}

\begin{theorem}[球上的 Poincaré 不等式]\label{theorem:ch05-球Poincare}
设 $1\leqslant p\leqslant\infty$。则存在只依赖于 $n$ 与 $p$的常数 $C$，使对每个球 $B(x,r)\subset\mathbb{R}^{n}$ 及每个函数 $u\in W^{1,p}(B^{0}(x,r))$ 有
\begin{align}
\|u-(u)_{x,r}\|_{L^{p}(B(x,r))}\leqslant Cr\|Du\|_{L^{p}(B(x,r))},
\end{align}
其中$(u)_{x,r}=\fint_{B(x,r)}u\,\mathrm{d}y$，即 $u$ 在球 $B(x,r)$ 上的平均.
\end{theorem}
\begin{remark}
如果直接在$B(x,r)$上使用\hyperref[theorem:ch05-Poincare]{Poincaré不等式}得到的常数$C$会与$B(x,r)$相关.
\end{remark}
\begin{proof}
$U=B^{0}(0,1)$ 的情形由\hyperref[theorem:ch05-Poincare]{Poincaré不等式}即得。一般地，若 $u\in W^{1,p}(B^{0}(x,r))$，写
\begin{align*}
v(y)\triangleq u(x+ry)\qquad (y\in B(0,1)).
\end{align*}
则 $v\in W^{1,p}(B^{0}(0,1))$，且
\begin{align*}
\|v-(v)_{0,1}\|_{L^{p}(B(0,1))}\leqslant C\|Dv\|_{L^{p}(B(0,1))}.
\end{align*}
作变量代换，再由\hyperref[theorem:ch05-Poincare]{Poincaré不等式}即得所述估计。
\end{proof}

\begin{definition}\label{definition:BMO}
设 $u\in L^1_{\mathrm{loc}}(\mathbb{R}^n)$。若
\begin{align*}
[u]_{BMO(\mathbb{R}^n)} \triangleq \sup\limits_{B(x,r)\subseteq \mathbb{R}^n} \left\{ \fint_{B(x,r)} |u - (u)_{x,r}| \mathrm{d}y \right\} < \infty,
\end{align*}
则称 $u$ 属于\textbf{有界平均振荡空间}$BMO(\mathbb{R}^n)$，并称 $[u]_{BMO(\mathbb{R}^n)}$ 为 $u$ 的\textbf{$BMO$半范数}。其中 $(u)_{x,r} = \fint_{B(x,r)} u \mathrm{d}y$ 表示 $u$ 在球 $B(x,r)$ 上的平均。
\end{definition}

\begin{theorem}\label{theorem:W1n嵌入BMO}
设 $u\in W^{1,n}(\mathbb{R}^n)\cap L^1(\mathbb{R}^n)$,则 $u\in BMO(\mathbb{R}^n)$，且存在只依赖于 $n$ 的常数 $C$，使得
\begin{align}\label{eq:BMO估计}
[u]_{BMO(\mathbb{R}^n)} \leqslant C \|Du\|_{L^n(\mathbb{R}^n)}.
\end{align}
\end{theorem}
\begin{proof}
任取球 $B(x,r)\subseteq \mathbb{R}^n$。由\hyperref[theorem:ch05-球Poincare]{球上的 Poincaré 不等式（$p=1$ 的情形）}得
\begin{align*}
\fint_{B(x,r)} |u - (u)_{x,r}| \mathrm{d}y \leqslant C r \fint_{B(x,r)} |Du| \mathrm{d}y.
\end{align*}
再由\hyperref[Real Analysis-theorem:实变函数--Hölder不等式]{Hölder不等式}得
\begin{align*}
\fint_{B(x,r)} |Du| \mathrm{d}y \leqslant \left( \fint_{B(x,r)} |Du|^n \mathrm{d}y \right)^{\frac{1}{n}}.
\end{align*}
注意到 $|B(x,r)| = \alpha(n) r^n$，其中 $\alpha(n)$ 为单位球体积，于是
\begin{align*}
r \left( \fint_{B(x,r)} |Du|^n \mathrm{d}y \right)^{\frac{1}{n}}
= r \left( \frac{1}{\alpha(n) r^n} \int_{B(x,r)} |Du|^n \mathrm{d}y \right)^{\frac{1}{n}}
= \frac{1}{(\alpha(n))^{\frac{1}{n}}} \left( \int_{B(x,r)} |Du|^n \mathrm{d}y \right)^{\frac{1}{n}}.
\end{align*}
因此
\begin{align*}
\fint_{B(x,r)} |u - (u)_{x,r}| \mathrm{d}y \leqslant \frac{C}{(\alpha(n))^{\frac{1}{n}}} \left( \int_{\mathbb{R}^n} |Du|^n \mathrm{d}y \right)^{\frac{1}{n}}.
\end{align*}
对所有的球 $B(x,r)\subseteq \mathbb{R}^n$ 取上确界，即得
\begin{align*}
[u]_{BMO(\mathbb{R}^n)} \leqslant C \|Du\|_{L^n(\mathbb{R}^n)}.
\end{align*}
故 $u\in BMO(\mathbb{R}^n)$。
\end{proof}



\subsection{几乎处处可微性}\label{subsection:ch05-a.e.可微}

\begin{definition}\label{definition:ch05-可微}
函数 $u:U\to\mathbb{R}$ 在 $x\in U$ 处可微，若存在 $a\in\mathbb{R}^{n}$ 使当 $y\to x$ 时
\begin{align}
u(y)=u(x)+a\cdot(y-x)+o(|y-x|).
\end{align}
换言之，有
\begin{align*}
\lim_{y\to x}\frac{|u(y)-u(x)-a\cdot(y-x)|}{|y-x|}=0.
\end{align*}
容易验证 $a$ 若存在则唯一。此后我们写 $Du(x)$ 表示 $a$，并称 $Du$ 为 $u$ 的\textbf{梯度}。
\end{definition}

\begin{theorem}[几乎处处可微]\label{theorem:ch05-a.e.可微}
设 $u\in W^{1,p}_{\mathrm{loc}}(U)$，其中 $n<p\leqslant\infty$。则 $u$ 在 $U$ 中几乎处处可微，且其梯度几乎处处等于其弱梯度。
\end{theorem}
\begin{remark}
回忆\cref{theorem:ch05-W1p有界估计},我们总把 $u$ 与其连续等价形等同.
\end{remark}
\begin{proof}

1. 先设 $n<p<\infty$。由\cref{corollary:ch05-Morrey局部估计}可得,对任何 $C^{1}$ 函数 $v$ 和$r>0$成立
\begin{align}\label{eq::jijiojiojiijji323399090ps}
|v(y)-v(x)|\leqslant Cr^{1-\frac{n}{p}}\left(\int_{B(x,2r)}|Dv(z)|^{p}\,\mathrm{d}z\right)^{\frac{1}{p}}\quad (y\in B(x,r)),
\end{align}
因而由逼近对任何 $v\in W^{1,p}$ 也成立。

2. 取 $u\in W^{1,p}_{\mathrm{loc}}(U)$。现对几乎处处的$x\in U$，由\hyperref[theorem:appE-Lebesgue微分1]{Lebesgue微分定理}可得，当$r\to 0$时,有
\begin{align}\label{eq::jijiojiojiijji323399090ps-1}
\fint_{B(x,r)}|Du(x)-Du(z)|^{p}\,\mathrm{d}z\to 0,
\end{align}
其中 $Du$ 如常表示 $u$ 的弱导数。固定任一这样的点 $x$，在估计中令
\begin{align*}
v(y)\triangleq u(y)-u(x)-Du(x)\cdot(y-x),
\end{align*}
其中$r=|x-y|.$当$r$充分小时,由\eqref{eq::jijiojiojiijji323399090ps}\eqref{eq::jijiojiojiijji323399090ps-1}式有
\begin{align*}
|u(y)-u(x)-Du(x)\cdot(y-x)|
&\leqslant Cr^{1-\frac{n}{p}}\left(\int_{B(x,2r)}|Du(x)-Du(z)|^{p}\,\mathrm{d}z\right)^{\frac{1}{p}}\\
&\leqslant Cr\left(\fint_{B(x,2r)}|Du(x)-Du(z)|^{p}\,\mathrm{d}z\right)^{\frac{1}{p}}\\
&=o(r)=o(|x-y|).
\end{align*}
故 $u$ 在 $x$ 处可微，且其梯度在 $x$ 处等于其弱梯度。

3. 若 $p=\infty$，即 $u\in W^{1,\infty}_{\mathrm{loc}}(U)$。取任意 $q$ 满足 $n<q<\infty$。我们首先证明 $u\in W^{1,q}_{\mathrm{loc}}(U)$。事实上，对任意紧集 $K\subset U$，由于 $u\in L^\infty(K)$ 且弱导数 $Du\in L^\infty(K)$，而 $K$ 的 Lebesgue 测度 $|K|<\infty$，故
\begin{align*}
\|u\|_{L^q(K)} \leqslant \|u\|_{L^\infty(K)} |K|^{\frac{1}{q}} < \infty,\qquad
\|Du\|_{L^q(K)} \leqslant \|Du\|_{L^\infty(K)} |K|^{\frac{1}{q}} < \infty.
\end{align*}
因此 $u\in W^{1,q}(K)$。由 $K$ 的任意性，$u\in W^{1,q}_{\mathrm{loc}}(U)$。于是对 $q$ 应用第 1、2 步的结论（注意 $n<q<\infty$），可知 $u$ 在 $U$ 中几乎处处可微，且其梯度几乎处处等于其弱梯度。由\hyperref[lemma:ch05-弱导数唯一性]{弱导数唯一性}知,弱梯度在 $W^{1,\infty}$ 与 $W^{1,q}$ 意义下一致，故 $p=\infty$ 时定理成立。
\end{proof}

\begin{theorem}[Rademacher 定理]\label{theorem:ch05-Rademacher}设 $u$ 在 $U$ 中局部 Lipschitz 连续。则 $u$ 在 $U$ 中几乎处处可微。
\end{theorem}
\begin{proof}
由\cref{theorem:ch05-W1infty局部刻画}和\cref{theorem:ch05-a.e.可微}的$p=\infty$情形立得.
\end{proof}

\subsection{Hardy 不等式}\label{subsection:ch05-Hardy}

\begin{theorem}[Hardy不等式]\label{theorem:ch05-Hardy}
设 $n\geqslant 3$ 且 $r>0$。设 $u\in H^1(B(0,r))$。则 $\frac{u}{|x|}\in L^2(B(0,r))$，且满足
\begin{align}\label{eq:Hardy-ineq-17}
\int_{B(0,r)} \frac{u^2}{|x|^2} \mathrm{d}x \leqslant 4 \int_{B(0,r)} \left( |Du|^2 + \frac{C u^2}{r^2} \right) \mathrm{d}x,
\end{align}
其中
\begin{align*}
C = \frac{n}{2(n-2 + \sqrt{(n-3)(n-1)})} \in \left(\frac{1}{4}, \frac{3}{2}\right].
\end{align*}
\end{theorem}
\begin{remark}
这不是\hyperref[theorem:Gagliardo--Nirenberg--Sobolev不等式]{Gagliardo--Nirenberg--Sobolev不等式}的推论.
\end{remark}
\begin{proof}
先设 $u\in C^\infty(B(0,r))$。注意到
\begin{align*}
D\left(\frac{1}{|x|}\right) = -\frac{x}{|x|^3}.
\end{align*}
由\hyperref[Basis of Analytics-theorem:Green恒等式]{Green恒等式\ref{Basis of Analytics-theorem:Green恒等式-3}}，对 $\bm{F} = \frac{u^2 x}{|x|}$ 和$g = \frac{1}{|x|}$，有
\begin{align*}
\int_{B(0,r)} \bm{F}\cdot \nabla g \,\mathrm{d}x
= \int_{\partial B(0,r)} g \bm{F}\cdot \nu \,\mathrm{d}S - \int_{B(0,r)} g \nabla\cdot \bm{F} \,\mathrm{d}x,
\end{align*}
其中$v=\frac{x}{|x|}$ 是 $\partial B(0,r)$ 的单位外法向量.
因为 $\nabla g = -\frac{x}{|x|^3}$，所以
\begin{align*}
\int_{B(0,r)} \bm{F}\cdot \nabla g \,\mathrm{d}x
= -\int_{B(0,r)} \frac{u^2 x}{|x|}\cdot \frac{x}{|x|^3} \,\mathrm{d}x
= -\int_{B(0,r)} \frac{u^2}{|x|^2} \,\mathrm{d}x.
\end{align*}
另一方面，直接计算得
\begin{align*}
\nabla\cdot \bm{F}
= \nabla\cdot \left( \frac{u^2 x}{|x|} \right)
= \frac{2u x\cdot Du}{|x|} + \frac{(n-1)u^2}{|x|}.
\end{align*}
在边界 $|x|=r$ 上，$\bm{F}\cdot \nu = u^2$，且 $g = \frac{1}{r}$，故
\begin{align*}
\int_{\partial B(0,r)} g \bm{F}\cdot \nu \,\mathrm{d}S
= \int_{\partial B(0,r)} \frac{u^2}{r} \,\mathrm{d}S.
\end{align*}
因此
\begin{align*}
-\int_{B(0,r)} \frac{u^2}{|x|^2} \,\mathrm{d}x
= \int_{\partial B(0,r)} \frac{u^2}{r} \,\mathrm{d}S
- \int_{B(0,r)} \frac{1}{|x|} \left( \frac{2u x\cdot Du}{|x|} + \frac{(n-1)u^2}{|x|} \right) \mathrm{d}x.
\end{align*}
整理得
\begin{align*}
\int_{B(0,r)} \frac{(n-2)u^2}{|x|^2} \,\mathrm{d}x
= -\int_{B(0,r)} \frac{2u x\cdot Du}{|x|^2} \,\mathrm{d}x
+ \int_{\partial B(0,r)} \frac{u^2}{r} \,\mathrm{d}S.
\end{align*}
应用\hyperref[Basis of Analytics-theorem:Gauss-Green公式(散度定理)]{Gauss-Green定理}于上式右边第二个积分，得
\begin{align}\label{eq:trace-18}
\int_{\partial B(0,r)} \frac{u^2}{r} \,\mathrm{d}S
= \int_{\partial B(0,r)} \frac{u^2 x}{r^2}\cdot \nu \,\mathrm{d}S
= \int_{B(0,r)} \nabla\cdot \left( \frac{u^2 x}{r^2} \right) \mathrm{d}x
= \int_{B(0,r)} \frac{n u^2 + 2u(x\cdot Du)}{r^2} \,\mathrm{d}x.
\end{align}
代入前式得
\begin{align*}
\int_{B(0,r)} \frac{(n-2)u^2}{|x|^2} \,\mathrm{d}x
= \int_{B(0,r)} \frac{n u^2}{r^2} \,\mathrm{d}x
+ \int_{B(0,r)} \frac{2u}{|x|^2} \left( \frac{|x|^2}{r^2} - 1 \right) x\cdot Du \,\mathrm{d}x.
\end{align*}
对上式右边第二个积分项应用\hyperref[theorem:带参的Cauchy不等式]{带参数 $\varepsilon$ 的 Cauchy 不等式}得
\begin{align*}
&\int_{B(0,r)}{\frac{2u}{|x|^2}\left( \frac{|x|^2}{r^2}-1 \right) x\cdot Du\,\mathrm{d}x}\leqslant \int_{B(0,r)}{\frac{2\left| u \right|}{|x|^2}\cdot \left| \frac{|x|^2}{r^2}-1 \right|\cdot \left| x \right|\cdot \left| Du \right|\,\mathrm{d}x}
\\
&\leqslant \int_{B(0,r)}{\frac{\left| u \right|}{|x|}\cdot 2\left| Du \right|\,\mathrm{d}x}\leqslant \varepsilon \int_{B(0,r)}{\frac{u^2}{|x|^2}\,\mathrm{d}x}+\frac{1}{\varepsilon}\int_{B(0,r)}{|Du|^2\,\mathrm{d}x}.
\end{align*}
于是
\begin{align*}
\int_{B(0,r)} \frac{(n-2)u^2}{|x|^2} \,\mathrm{d}x
\leqslant \int_{B(0,r)} \frac{n u^2}{r^2} \,\mathrm{d}x
+ \varepsilon \int_{B(0,r)} \frac{u^2}{|x|^2} \,\mathrm{d}x
+ \frac{1}{\varepsilon} \int_{B(0,r)} |Du|^2 \,\mathrm{d}x.
\end{align*}
移项得
\begin{align*}
\int_{B(0,r)} \frac{(n-2-\varepsilon)u^2}{|x|^2} \,\mathrm{d}x
\leqslant \frac{1}{\varepsilon} \int_{B(0,r)} |Du|^2 \,\mathrm{d}x
+ \int_{B(0,r)} \frac{n u^2}{r^2} \,\mathrm{d}x.
\end{align*}
因 $n\geqslant 3$，令 $\varepsilon(n-2-\varepsilon) = \lambda > 0$，只要
\begin{align*}
\varepsilon = \frac{n-2 - \sqrt{(n-2)^2 - 4\lambda}}{2} > 0
\end{align*}
即可。
\begin{align*}
\int_{B(0,r)} \frac{u^2}{|x|^2} \,\mathrm{d}x
\leqslant \frac{1}{\lambda} \left( \int_{B(0,r)} |Du|^2 \,\mathrm{d}x
+ n\varepsilon \int_{B(0,r)} \frac{u^2}{r^2} \,\mathrm{d}x \right)
.
\end{align*}
固定$n$,$\lambda=\varepsilon(n-2-\varepsilon)$是关于$\varepsilon$的二元函数,最大值为$\lambda_{\text{max}}=\frac{(n-2)^2}{4}$.因为不等式要对所有$n\geqslant 3$成立,所以$\lambda$最大只可取 $\lambda = \frac{1}{4}$，此时令
$4n\varepsilon = 4C$，则
\begin{align*}
C = \frac{n}{2(n-2 + \sqrt{(n-3)(n-1)})}.
\end{align*}
再代入上式得
\begin{align*}
\int_{B(0,r)} \frac{u^2}{|x|^2} \,\mathrm{d}x
\leqslant 4 \int_{B(0,r)} |Du|^2 \,\mathrm{d}x
+ 4n\varepsilon \int_{B(0,r)} \frac{u^2}{r^2} \,\mathrm{d}x= 4 \int_{B(0,r)} \left( |Du|^2 + \frac{C u^2}{r^2} \right) \mathrm{d}x.
\end{align*}
最后，对一般 $u\in H^1(B(0,r))$，由\hyperref[theorem:ch05-边界光滑逼近]{边界光滑函数的整体逼近定理},可取 $u_m\in C^\infty(\overline{B(0,r)})$ 使得 $u_m\to u$ 于 $H^1$。对每个 $u_m$ 上述不等式成立。由\hyperref[Real Analysis-theorem:Fatou引理]{Fatou引理}得
\begin{align*}
\int_{B(0,r)} \frac{u^2}{|x|^2} \,\mathrm{d}x
\leqslant \varliminf_{m\to\infty} \int_{B(0,r)} \frac{u_m^2}{|x|^2} \,\mathrm{d}x \leqslant \lim_{m\to\infty} 4 \int_{B(0,r)} \left( |Du_m|^2 + \frac{C u_m^2}{r^2} \right) \mathrm{d}x
= 4 \int_{B(0,r)} \left( |Du|^2 + \frac{C u^2}{r^2} \right) \mathrm{d}x.
\end{align*}
\end{proof}
\begin{remark}
应用\hyperref[theorem:带参的Cauchy不等式]{带参数 $\varepsilon$ 的 Cauchy 不等式}于 \eqref{eq:trace-18} 式右边得迹不等式
\begin{align*}
\frac{1}{r} \int_{\partial B(0,r)} u^2 \mathrm{d}S
\leqslant C \int_{B(0,r)} \left( |Du|^2 + \frac{u^2}{r^2} \right) \mathrm{d}x,
\end{align*}
其中
\begin{align*}
C = \frac{n + \sqrt{n^2+4}}{2}.
\end{align*}
它不限于 $n\geqslant 3$。习题 7 是另一种得到迹不等式的方法。习题 16 是 \eqref{eq:Hardy-ineq-17} 中 $r=\infty$ 时的变形的 Hardy 不等式。
\end{remark}


\subsection{Fourier变换方法}\label{subsection:ch05-Fourier}

下面我们利用Fourier变换（\cref{section:ch04-变换方法}）给出空间 $H^{k}(\mathbb{R}^{n})$ 的另一种刻画。本小节所有函数均为复值。

\begin{theorem}[用 Fourier 变换刻画 $H^k$]\label{theorem:Fourier刻画Hk}
设 $k$ 为非负整数。
函数 $u\in L^2(\mathbb{R}^n)$ 属于 $H^k(\mathbb{R}^n)$ 当且仅当
\begin{align}\label{eq:Fourier-Hk-19}
(1+|y|^k)\hat{u}\in L^2(\mathbb{R}^n).
\end{align}
此外, 对 $\forall u\in H^k(\mathbb{R}^n)$,
\begin{align}\label{eq:Fourier-Hk-20}
\frac{1}{\sqrt{\max\{k,1\}}}\|u\|_{H^k(\mathbb{R}^n)}\leqslant \|(1+|y|^k)\hat{u}\|_{L^2(\mathbb{R}^n)}\leqslant \sqrt{2}\|u\|_{H^k(\mathbb{R}^n)}.
\end{align}
\end{theorem}
\begin{proof}

{\heiti 1. } 当 $k=0$ 时定理自然成立. 现设 $k\geqslant 1$. 首先, 假设 $u\in H^k(\mathbb{R}^n)$. 则由 $H^k(\mathbb{R}^n)$ 的定义, 对每一多重指标 $|\alpha|\leqslant k$, 我们有 $D^\alpha u\in L^2(\mathbb{R}^n)$. 若 $u\in C_c^k(\mathbb{R}^n)$ 有紧支集, 根据\hyperref[proposition:Fourier 变换的基本性质2]{Fourier变换性质\ref{proposition:Fourier 变换的基本性质2-2}}得
\begin{align}\label{eq:Fourier-Hk-21}
(\text{i}y)^\alpha\hat{u}=\widehat{D^\alpha u}.
\end{align}
当 $u\in H^k(\mathbb{R}^n)$ 时, 由\hyperref[theorem:ch05-边界光滑逼近]{边界光滑函数的整体逼近定理},取一列光滑函数 $u_m\in C_c^\infty(\mathbb{R}^n)$ 使得 $u_m\to u$ 于 $H^k(\mathbb{R}^n)$, 并用光滑函数逼近可推出公式\eqref{eq:Fourier-Hk-21} 对一般的 $u\in H^k(\mathbb{R}^n)$ 也成立. 故对每个 $|\alpha|\leqslant k$, $(\text{i}y)^\alpha\hat{u}\in L^2(\mathbb{R}^n)$. 由\eqref{eq:Fourier-Hk-21}式及\hyperref[theorem:Plancherel定理]{Plancherel定理}得
\begin{align*}
\int_{\mathbb{R} ^n}{|y|^{2k}|\hat{u}|^2\,\mathrm{d}y}&\overset{?}{=}\int_{\mathbb{R} ^n}{\sum_{|\alpha |=k}{\left| y^{\alpha} \right|^2\left| \hat{u} \right|^2\,\mathrm{d}x}}=\int_{\mathbb{R} ^n}{\sum_{|\alpha |=k}{\left| \left( \mathrm{i}y \right) ^{\alpha}\hat{u} \right|^2\,\mathrm{d}x}}
\\
&=\int_{\mathbb{R} ^n}{\sum_{|\alpha |=k}{|D^{\alpha}u|^2\,\mathrm{d}x}}=\int_{\mathbb{R} ^n}{|D^ku|^2\,\mathrm{d}x}<\infty .
\end{align*}
所以, 因 $(1+|y|^k)^2\leqslant 2(1+|y|^{2k})$, 我们有
\begin{align*}
\int_{\mathbb{R}^n}(1+|y|^k)^2|\hat{u}|^2\,\mathrm{d}y
\leqslant 2\int_{\mathbb{R}^n}(1+|y|^{2k})|\hat{u}|^2\,\mathrm{d}y
\leqslant 2\|u\|_{H^k(\mathbb{R}^n)}^2<\infty.
\end{align*}
从而 $(1+|y|^k)\hat{u}\in L^2(\mathbb{R}^n)$. 同时由上式还可得
\begin{align*}
\|(1+|y|^k)\hat{u}\|_{L^2(\mathbb{R}^n)}\leqslant \sqrt{2}\|u\|_{H^k(\mathbb{R}^n)}.
\end{align*}
这就是 \eqref{eq:Fourier-Hk-20} 式的右半边不等式.

{\heiti 2. } 反之, 若 $(1+|y|^k)\hat{u}\in L^2(\mathbb{R}^n)$, 且 $|\alpha|\leqslant k$, 则当 $k\geqslant 1$ 时,有
\begin{align}\label{eq:Fourier-Hk-22}
\sum_{s=0}^k\int_{\mathbb{R}^n}|y|^{2s}|\hat{u}|^2\,\mathrm{d}y
\leqslant k\int_{\mathbb{R}^n}(1+|y|^k)^2|\hat{u}|^2\,\mathrm{d}y.
\end{align}
令
\begin{align*}
u_\alpha\triangleq ((\text{i}y)^\alpha\hat{u})^\vee.
\end{align*}
则对每一 $\phi\in C_c^\infty(\mathbb{R}^n)$, 由\hyperref[theorem:Plancherel定理]{Plancherel定理}及\hyperref[proposition:Fourier 变换的基本性质2]{Fourier变换性质\ref{proposition:Fourier 变换的基本性质2-2}}得
\begin{align*}
\int_{\mathbb{R}^n}(D^\alpha\phi)\bar{u}\,\mathrm{d}x
=\int_{\mathbb{R}^n}\widehat{(D^\alpha\phi)}\bar{\hat{u}}\,\mathrm{d}y
=\int_{\mathbb{R}^n}(\text{i}y)^\alpha\hat{\phi}\bar{\hat{u}}\,\mathrm{d}y
=(-1)^{|\alpha|}\int_{\mathbb{R}^n}\hat{\phi}\overline{(\text{i}y)^\alpha\hat{u}}\,\mathrm{d}y
=(-1)^{|\alpha|}\int_{\mathbb{R}^n}\phi\bar{u}_\alpha\,\mathrm{d}x.
\end{align*}
因而在弱意义下 $u_\alpha=D^\alpha u$, 所以
\begin{align*}
\sum_{s=0}^k\int_{\mathbb{R}^n}|D^su|^2\,\mathrm{d}x
=\sum_{s=0}^k\int_{\mathbb{R}^n}|y|^{2s}|\hat{u}|^2\,\mathrm{d}y.
\end{align*}
由 \eqref{eq:Fourier-Hk-22} 得 $u\in H^k(\mathbb{R}^n)$, 且由上式与 \eqref{eq:Fourier-Hk-22} 有
\begin{align*}
\|u\|_{H^k(\mathbb{R}^n)}^2
=\sum_{|\alpha|\leqslant k}\int_{\mathbb{R}^n}|D^\alpha u|^2\,\mathrm{d}x
\leqslant \max\{k,1\}\int_{\mathbb{R}^n}(1+|y|^k)^2|\hat{u}|^2\,\mathrm{d}y.
\end{align*}
故有 \eqref{eq:Fourier-Hk-20} 式的左半边估计式. 
\end{proof}

\begin{lemma}\label{lemma:多项式等价}
设 $k$ 为非负整数。则存在只依赖于 $n$ 与 $k$ 的常数 $c_1,c_2>0$，使得对一切 $y\in\mathbb{R}^n$，
\begin{align}\label{eq:多项式等价}
c_1(1+|y|^k)^2
\leqslant \sum_{|\alpha|\leqslant k}|y^\alpha|^2
\leqslant c_2(1+|y|^k)^2.
\end{align}
\end{lemma}
\begin{proof}
记
\begin{align*}
P(y)\triangleq\sum_{|\alpha|\leqslant k}|y^\alpha|^2,\qquad Q(y)\triangleq(1+|y|^k)^2.
\end{align*}
则 $P,Q$ 都是连续函数，且 $P(y)\geqslant 1$，$Q(y)\geqslant 1$。

{\heiti 上界. } 对每个 $|\alpha|\leqslant k$，有 $|y^\alpha|\leqslant (1+|y|)^{|\alpha|}\leqslant (1+|y|)^k$。因此
\begin{align*}
P(y)\leqslant \sum_{|\alpha|\leqslant k}(1+|y|)^{2k}
\leqslant N_{n,k}(1+|y|)^{2k},
\end{align*}
其中 $N_{n,k}$ 表示满足 $|\alpha|\leqslant k$ 的多重指标个数，只依赖于 $n,k$。又因为
\begin{align*}
(1+|y|)^{2k}\leqslant 2^{2k}(1+|y|^{2k})\leqslant 2^{2k}(1+|y|^k)^2,
\end{align*}
故取 $c_2\triangleq 2^{2k}N_{n,k}$，即得
\begin{align*}
P(y)\leqslant c_2(1+|y|^k)^2.
\end{align*}

{\heiti 下界. } 先证存在常数 $A>0$ 使得
\begin{align}\label{eq:齐次控制}
|y|^{2k}\leqslant A\sum_{|\alpha|=k}|y^\alpha|^2.
\end{align}
事实上，由多项式展开，
\begin{align*}
|y|^{2k}=\left(\sum_{i=1}^n y_i^2\right)^k
=\sum_{|\alpha|=k}\frac{k!}{\alpha!}y^{2\alpha}
=\sum_{|\alpha|=k}\frac{k!}{\alpha!}|y^\alpha|^2
\leqslant k!\sum_{|\alpha|=k}|y^\alpha|^2.
\end{align*}
故可取 $A=k!$。于是
\begin{align*}
|y|^k\leqslant \sqrt{k!}\left(\sum_{|\alpha|=k}|y^\alpha|^2\right)^{1/2}.
\end{align*}
令 $X\triangleq\sum_{|\alpha|=k}|y^\alpha|^2\geqslant 0$，则
\begin{align*}
Q(y)&=(1+|y|^k)^2
=1+2|y|^k+|y|^{2k}
\leqslant 1+2\sqrt{k!}\sqrt{X}+k!X.
\end{align*}
因为 $2\sqrt{X}\leqslant 1+X$，所以
\begin{align*}
1+2\sqrt{k!}\sqrt{X}+k!X
&\leqslant 1+\sqrt{k!}(1+X)+k!X\\
&=(1+\sqrt{k!})+( \sqrt{k!}+k!)X\\
&\leqslant (1+\sqrt{k!}+k!)(1+X).
\end{align*}
又因为
\begin{align*}
P(y)=\sum_{|\alpha|\leqslant k}|y^\alpha|^2
\geqslant 1+\sum_{|\alpha|=k}|y^\alpha|^2=1+X,
\end{align*}
故
\begin{align*}
Q(y)\leqslant (1+\sqrt{k!}+k!)P(y).
\end{align*}
取 $c_1\triangleq \frac{1}{1+\sqrt{k!}+k!}$，即得
\begin{align*}
c_1(1+|y|^k)^2\leqslant \sum_{|\alpha|\leqslant k}|y^\alpha|^2.
\end{align*}
\end{proof}

\begin{theorem}[用Fourier变换刻画 $H^k$]\label{theorem:Fourier刻画Hk-严谨}
设 $k$ 为非负整数。则 $u\in H^k(\mathbb{R}^n)$ 当且仅当
\begin{align*}
(1+|y|^k)\hat{u}\in L^2(\mathbb{R}^n).
\end{align*}
此外，存在只依赖于 $n$ 与 $k$ 的常数 $C_1,C_2>0$，使得对一切 $u\in H^k(\mathbb{R}^n)$，
\begin{align}\label{eq:Fourier-Hk-等价范数11}
C_1\|u\|_{H^k(\mathbb{R}^n)}
\leqslant \|(1+|y|^k)\hat{u}\|_{L^2(\mathbb{R}^n)}
\leqslant C_2\|u\|_{H^k(\mathbb{R}^n)}.
\end{align}
\end{theorem}
\begin{proof}

{\heiti 1. } 当 $k=0$ 时定理自然成立. 现设 $k\geqslant 1$. 首先, 假设 $u\in H^k(\mathbb{R}^n)$. 则由 $H^k(\mathbb{R}^n)$ 的定义, 对每一多重指标 $|\alpha|\leqslant k$, 我们有 $D^\alpha u\in L^2(\mathbb{R}^n)$. 若 $u\in C_c^k(\mathbb{R}^n)$ 有紧支集, 根据\hyperref[theorem:1222-Fourier变换的性质]{Fourier变换性质\ref{theorem:1222-Fourier变换的性质-1}}得
\begin{align}\label{eq:Fourier-Hk-211}
(\text{i}y)^\alpha\hat{u}=\widehat{D^\alpha u}.
\end{align}
当 $u\in H^k(\mathbb{R}^n)$ 时, 由\hyperref[theorem:ch05-边界光滑逼近]{边界光滑函数的整体逼近定理}, 取一列光滑函数 $u_m\in C_c^\infty(\mathbb{R}^n)$ 使得 $u_m\to u$ 于 $H^k(\mathbb{R}^n)$, 并用光滑函数逼近可推出公式 \eqref{eq:Fourier-Hk-211} 对一般的 $u\in H^k(\mathbb{R}^n)$ 也成立. 故对每个 $|\alpha|\leqslant k$, $(\text{i}y)^\alpha\hat{u}\in L^2(\mathbb{R}^n)$. 由 \eqref{eq:Fourier-Hk-211} 式及\hyperref[theorem:Plancherel定理]{Plancherel定理}得
\begin{align*}
\int_{\mathbb{R}^n}\sum_{|\alpha|=k}|y^\alpha|^2|\hat{u}|^2\,\mathrm{d}y
&=\int_{\mathbb{R}^n}\sum_{|\alpha|=k}|(\text{i}y)^\alpha\hat{u}|^2\,\mathrm{d}y
=\int_{\mathbb{R}^n}\sum_{|\alpha|=k}|\widehat{D^\alpha u}|^2\,\mathrm{d}y \\
&=\int_{\mathbb{R}^n}\sum_{|\alpha|=k}|D^\alpha u|^2\,\mathrm{d}x
=\int_{\mathbb{R}^n}|D^k u|^2\,\mathrm{d}x<\infty.
\end{align*}
因此, 由\cref{lemma:多项式等价}, 存在只依赖于 $n$ 与 $k$ 的常数 $c_1>0$, 使得
\begin{align*}
(1+|y|^k)^2\leqslant c_1\sum_{|\alpha|\leqslant k}|y^\alpha|^2.
\end{align*}
于是
\begin{align*}
\int_{\mathbb{R}^n}(1+|y|^k)^2|\hat{u}|^2\,\mathrm{d}y
\leqslant c_1\int_{\mathbb{R}^n}\sum_{|\alpha|\leqslant k}|y^\alpha|^2|\hat{u}|^2\,\mathrm{d}y
\leqslant c_1\|u\|_{H^k(\mathbb{R}^n)}^2<\infty.
\end{align*}
从而 $(1+|y|^k)\hat{u}\in L^2(\mathbb{R}^n)$. 同时由上式还可得
\begin{align*}
\|(1+|y|^k)\hat{u}\|_{L^2(\mathbb{R}^n)}\leqslant C_2\|u\|_{H^k(\mathbb{R}^n)},
\end{align*}
其中 $C_2=\sqrt{c_1}$. 这就是 \eqref{eq:Fourier-Hk-等价范数11} 式的右半边不等式.

{\heiti 2. } 反之, 若 $(1+|y|^k)\hat{u}\in L^2(\mathbb{R}^n)$, 且 $|\alpha|\leqslant k$, 则当 $k\geqslant 1$ 时, 由\cref{lemma:多项式等价}, 存在只依赖于 $n$ 与 $k$ 的常数 $c_2>0$, 使得
\begin{align}\label{eq:Fourier-Hk-222}
\sum_{|\alpha|\leqslant k}|y^\alpha|^2
\leqslant c_2(1+|y|^k)^2.
\end{align}
令
\begin{align*}
u_\alpha\triangleq ((\text{i}y)^\alpha\hat{u})^\vee.
\end{align*}
则对每一 $\phi\in C_c^\infty(\mathbb{R}^n)$, 由\hyperref[theorem:Plancherel定理]{Plancherel定理}及\hyperref[theorem:1222-Fourier变换的性质]{Fourier变换性质\ref{theorem:1222-Fourier变换的性质-1}\ref{theorem:1222-Fourier变换的性质-2}}得
\begin{align*}
\int_{\mathbb{R}^n}(D^\alpha\phi)\bar{u}\,\mathrm{d}x
=\int_{\mathbb{R}^n}\widehat{(D^\alpha\phi)}\bar{\hat{u}}\,\mathrm{d}y
=\int_{\mathbb{R}^n}(\text{i}y)^\alpha\hat{\phi}\bar{\hat{u}}\,\mathrm{d}y =(-1)^{|\alpha|}\int_{\mathbb{R}^n}\hat{\phi}\overline{(\text{i}y)^\alpha\hat{u}}\,\mathrm{d}y
=(-1)^{|\alpha|}\int_{\mathbb{R}^n}\phi\bar{u}_\alpha\,\mathrm{d}x.
\end{align*}
因而在弱意义下 $u_\alpha=D^\alpha u$, 所以
\begin{align*}
\sum_{|\alpha|\leqslant k}\int_{\mathbb{R}^n}|D^\alpha u|^2\,\mathrm{d}x
=\sum_{|\alpha|\leqslant k}\int_{\mathbb{R}^n}|y^\alpha|^2|\hat{u}|^2\,\mathrm{d}y.
\end{align*}
由 \eqref{eq:Fourier-Hk-222} 得 $u\in H^k(\mathbb{R}^n)$, 且由上式与 \eqref{eq:Fourier-Hk-222} 有
\begin{align*}
\|u\|_{H^k(\mathbb{R}^n)}^2
=\sum_{|\alpha|\leqslant k}\int_{\mathbb{R}^n}|D^\alpha u|^2\,\mathrm{d}x
\leqslant c_2\int_{\mathbb{R}^n}(1+|y|^k)^2|\hat{u}|^2\,\mathrm{d}y.
\end{align*}
故有 \eqref{eq:Fourier-Hk-等价范数11} 式的左半边估计式, 其中 $C_1=\frac{1}{\sqrt{c_2}}$. 证毕.
\end{proof}

\begin{definition}\label{definition:ch05-Hs}
设 $0<s<\infty$，$u\in L^{2}(\mathbb{R}^{n})$。则 $u\in H^{s}(\mathbb{R}^{n})$ 若 $(1+|y|^{s})\widehat{u}\in L^{2}(\mathbb{R}^{n})$。对非整数 $s$，令
\begin{align*}
\|u\|_{H^{s}(\mathbb{R}^{n})}\triangleq\|(1+|y|^{s})\widehat{u}\|_{L^{2}(\mathbb{R}^{n})}.
\end{align*}
\end{definition}

\end{document}
